\begin{remark}
    Throughout this textbook, \textbf{logical formulas} are interpreted according to the following conventions:
    \begin{itemize}
        \item \textbf{Parentheses} have the highest precedence. For example,
        \[
        (P \lor Q) \land R
        \]
        is evaluated by first computing $P \lor Q$
        \item The \textbf{logical connectives} are interpreted in the order
        \[
        \neg \; > \; \land \; > \; \lor \; > \; \implies \; > \; \iff
        \]
        Hence
        \[
        \neg P \land Q \lor R \implies S \iff T
        \]
        is understood as
        \[
        ((((\neg P) \land Q) \lor R) \implies S) \iff T
        \]
        \item A \textbf{quantifier} binds the entire formula that follows it
        unless parentheses indicate a smaller scope. Thus
        \[
        \forall x, P(x) \land Q(x) \implies R(x)
        \]
        means
        \[
        \forall x\, ((P(x)\land Q(x))\implies R(x))
        \]
        whereas
        \[
        (\forall x P(x)) \land (\forall x Q(x))
        \]
        is different from
        \[
        \forall x (P(x) \land Q(x))
        \]
    \end{itemize}
\end{remark}

\begin{section}{The Natural Number Systems}

The set $\set*{1, 2, 3, \dots}$ of natural numbers is denoted by $\N$, and is defined by the following axioms:
\begin{itemize}
    \item $1 \in \N$
    \item $\forall n \in \N, s(n) \in \N$
    \item $\nexists n \in \N, s(n) = 1$
    \item $\forall m, n \in \N, s(m) = s(n) \implies m = n$
    \item If $S \subseteq \N$ satisfies
    \begin{itemize}
        \item $1 \in S$
        \item $\forall n \in S, s(n) \in S$
    \end{itemize}
    then $S = \N$
\end{itemize}
where $s(n)$ is called the \textbf{successor} of $n$.
These axioms are called the \textbf{Peano axioms} (or \emph{Peano Postulates}).

\bigskip

\begin{remark}
    The \textbf{Peano axioms} do not guarantee the \emph{existence} of $\N$, which is guaranteed by the
    \textbf{axiom of infinity} in the \textbf{ZFC axioms}. In this textbook, the \textbf{ZFC axioms}
    is not discussed, so we will just assume that $\N$ exists.
\end{remark}

\bigskip

\begin{definition} \label{def:MathematicalInduction}
    Let $P(n)$ be a statement about $n \in \N$, and let
    \[
    S \coloneqq \setbuilder*{n \in \N}{P(n) \text{ is true}} \subseteq \N
    \]
    Assume that $P(n)$ satisfies
    \begin{itemize}
        \item $P(1)$ is true
        \item $P(n)$ is true $\implies P(n + 1)$ is true
    \end{itemize}
    This implies that
    \begin{itemize}
        \item $1 \in S$
        \item $n \in S \implies n + 1 \in S$
    \end{itemize}
    Therefore,
    \[
    \begin{array}{c}
        S = \N \\
        \Downarrow \\
        P(n) \text{ is true for all } n \in \N
    \end{array}
    \]
    This method of proof is called the \textbf{principle of mathematical induction}.
\end{definition}

\bigskip

\begin{definition} \label{def:OperationsOnN}
    For all $n, m \in \N$, we define the \textbf{addition} operator $+$ on $\N$ by
    \begin{itemize}
        \item $n + 1 \coloneqq s(n)$
        \item $n + s(m) \coloneqq s(n + m)$
    \end{itemize}
    For all $n, m \in \N$, we define the \textbf{multiplication} operator $\cdot$ on $\N$ by
    \begin{itemize}
        \item $n \cdot 1 \coloneqq n$
        \item $n \cdot s(m) \coloneqq n \cdot m + n$
    \end{itemize}
\end{definition}

\bigskip

\begin{remark}
    Strictly speaking, \Cref{def:OperationsOnN} requires the \textbf{recursion theorem} to be strictly defined.
    In this textbook, the \textbf{recursion theorem} is not discussed, so we will just assume that the operators
    on $\N$ are well-defined.
    \\[6mm]
    Also, we will not discuss the properties of the operators on $\N$, such as
    \begin{itemize}
        \item $n + m = m + n$
        \item $(n + m) + k = n + (m + k)$
        \item $n \cdot m = m \cdot n$
        \item $(n \cdot m) \cdot k = n \cdot (m \cdot k)$
        \item $n \cdot (m + k) = n \cdot m + n \cdot k$
    \end{itemize}
    for all $n, m, k \in \N$. For these proofs rely on the recursive definitions.
\end{remark}

\bigskip

\begin{theorem} \label{thm:WellOrderingProperty}
    Let $S \subseteq \N$ be nonempty. Then $S$ has a least element.
    This theorem is called the \textbf{well-ordering property} of $\N$.
\end{theorem}

\begin{proof}
    Assume that $S$ has no least element. Let
    \begin{itemize}
        \item $T \coloneqq \N \setminus S$
        \item $L \coloneqq \setbuilder*{n \in \N}{\set*{1, 2, \dots, n} \subseteq T}$
    \end{itemize}
    We will show that
    \[
    \begin{array}{c}
        L = \N \\
        \Downarrow \\
        T = \N \\
        \Downarrow \\
        S = \vnot \\
        \Downarrow \\
        \bot
    \end{array}
    \]
    \begin{itemize}
        \item \textbf{Base case:} If $1 \in S$, then
        \[
        \begin{array}{c}
            1 \in S \\
            \Downarrow \\
            1 = \min S \\
            \Downarrow \\
            \bot
        \end{array}
        \]
        Therefore,
        \[
        \begin{array}{c}
            1 \notin S \\
            \Downarrow \\
            1 \in T \\
            \Downarrow \\
            1 \in L
        \end{array}
        \]
        \item \textbf{Inductive step:} Assume that $n \in L$ for some $n \in \N$. Then
        \[
        \begin{array}{c}
            n \in L \\
            \Downarrow \\
            \set*{1, 2, \dots, n} \subseteq T \\
            \Downarrow \\
            \set*{1, 2, \dots, n} \cap S = \vnot
        \end{array}
        \]
        If $n + 1 \in S$, then 
        \[
        \begin{array}{c}
            n + 1 \in S \\
            \Downarrow \\
            n + 1 = \min S \\
            \Downarrow \\
            \bot
        \end{array}
        \]
        Therefore,
        \[
        \begin{array}{c}
            n + 1 \notin S \\
            \Downarrow \\
            n + 1 \in T \\
            \Downarrow \\
            \set*{1, 2, \dots, n + 1} \subseteq T \\
            \Downarrow \\
            n + 1 \in L
        \end{array}
        \]
    \end{itemize}
    By the \emph{principle of mathematical induction}, $L = \N$.
\end{proof}

\end{section}